% =============================================================================
%  RAG SYSTEM — detailed methods (included from methodology.tex)
% =============================================================================

\subsection{RAG System}
\label{sec:rag-system}

This section presents the RAG system, which is the central technical
contribution of this thesis. The pipeline has six components: data
acquisition, digitization, database construction, embedding model selection,
two-stage retrieval, and answer generation.


% ─────────────────────────────────────────────────────────────────────────────
\subsubsection{Data Acquisition}
\label{subsec:data-acquisition}

The foundation of any RAG system is the quality of its knowledge
base~\cite{lewis2020rag, gao2024rag_survey}. Research on data quality in
AI systems has shown that downstream performance is often dominated by data
quality rather than model architecture, and that compounding quality issues
from early data stages are the most common cause of
failure~\cite{sambasivan2021everyone}. This guided our approach: we
prioritized corpus quality and structural consistency over size.

No publicly available, machine-readable corpus of Tunisian Bac mathematics
materials existed before this work. The primary sources are printed
examination booklets, handwritten corrections, and scanned course pages
from preparatory-class teachers.

\paragraph{Physical data collection.}
The raw material was collected during a trip to Tunisia, where we obtained
books and course materials from experienced Bac mathematics teachers. For
printed material, a standard scanner or phone camera produced sufficient
quality. For handwritten corrections with dense notation or faded ink, the
corrections were manually rewritten on a tablet (iPad with Apple Pencil)
before scanning---a labor-intensive but necessary investment in data
quality~\cite{sambasivan2021everyone, wang2023ocrbench,
mahdavi2019icdar_math}.

Each exercise was scanned \textit{together with its correction} as a single
image or PDF, rather than separating statements from solutions. Course
material was scanned as continuous documents per chapter.

\paragraph{Corpus structure.}
The material was uploaded to Google Cloud Storage~\cite{gcloud_storage2024}
in a hierarchical structure:

\begin{verbatim}
BacMath_Raw_Data/
  01_Continuite_et_limites/
    Cours/                      -> course material
    series_et_corrections/      -> exercises + corrections
  ...
  20_Coniques/
  Bac_avec_corrections_Images/  -> past exam corrections
\end{verbatim}

This naming convention is machine-readable: the indexing pipeline extracts
document type, chapter, year, session, and exercise number directly from
folder paths using regular expressions. The corpus comprises 624 source
files spanning 20 chapters: 312 corrected exam pages, 250 series exercise
files, and 62 course material files, yielding 1{,}422 chunks after indexing.

\paragraph{Paired exercise--correction acquisition.}
Scanning each exercise together with its correction was a deliberate design
decision. In mathematical pedagogy, the correction is a complete reasoning
chain whose structure depends on the exercise statement. Separating them
would sever this dependency. The retrieval system also benefits: the most
useful retrieval target is a complete worked example, not an orphaned
solution. This aligns with research showing that structure-aware document
processing improves retrieval
quality~\cite{zhao2024financial_chunking, chen2024mog}.


% ─────────────────────────────────────────────────────────────────────────────
\subsubsection{Mathematical Document Digitization}
\label{subsec:digitization}

The next step converts scanned images into machine-readable \LaTeX{}.
Mathematical documents contain two-dimensional structures (fractions,
subscripts, integrals) that standard OCR cannot
handle~\cite{wang2023ocrbench, mahdavi2019icdar_math}.

\paragraph{Tool selection.}
We initially considered Mathpix~\cite{mathpix2024}, a commercial math-OCR
service, but its per-page pricing made it impractical for a thesis prototype.
Instead, we used Gemini via Vertex AI. Benchmarking by Wang et
al.~\cite{wang2023ocrbench} found that Gemini achieved the highest overall
OCR score among multimodal models tested. For our material---primarily
printed or cleanly rewritten---Gemini provided sufficient fidelity at
negligible cost. The choice was driven by budget constraints, not a claim of
superior accuracy over Mathpix.

\paragraph{Pipeline (\texttt{digitize.py}).}
The digitization pipeline: (1)~enumerates all images and PDFs in GCS,
(2)~checks for existing companion \texttt{.tex} files (making the pipeline
resume-friendly), (3)~sends each pending file to Gemini with a structured
prompt requesting faithful \LaTeX{} transcription in the official redaction
style, and (4)~uploads the result back to GCS. The pipeline includes
exponential-backoff retry logic for transient API errors.

\paragraph{Quality.}
Digitization errors propagate through the entire system: a misrecognized
formula gets embedded, retrieved, and injected into the LLM's prompt. The
ground truth evaluation (Section~\ref{subsec:ocr-evaluation}) confirmed an
average CER of 5.4\% and \LaTeX{} math accuracy of 97.1\%---sufficient for
our corpus but not error-free.


% ─────────────────────────────────────────────────────────────────────────────
\subsubsection{Database Construction and Indexing}
\label{subsec:database-construction}

\paragraph{Migration to ChromaDB.}
An early prototype used a JSON-based vector store. This was functional for
initial experiments but could not scale (full memory load per query, no
incremental updates, no metadata filtering). We migrated to
ChromaDB~\cite{chromadb2024}, which provides persistent local storage,
native metadata filtering, and deterministic upsert operations.

\paragraph{Text normalization.}
Before chunking, each \texttt{.tex} file is normalized: CRLF $\to$ LF,
collapsed whitespace, and excess newlines reduced to double newlines.

\paragraph{Adaptive chunking.}
Corrections use 1{,}500-character chunks with 200-character overlap to
preserve exercise boundaries. Course material uses 3{,}000-character chunks
with 250-character overlap to keep theorems intact. The algorithm prefers
natural break points (paragraph breaks, line breaks, sentence endings) over
arbitrary character positions, consistent with findings that structure-aware
chunking improves retrieval~\cite{zhao2024financial_chunking, chen2024mog}.

\paragraph{Metadata extraction.}
Each chunk is annotated with metadata extracted from the GCS path:
\texttt{type} (bac\_officiel, serie, cours), \texttt{chapter},
\texttt{year}, \texttt{session}, \texttt{exo\_id}, \texttt{is\_solution}
(boolean), \texttt{group\_id} (for companion pairing), \texttt{filename},
and \texttt{source} (GCS URI). This metadata enables filtered retrieval and
exercise--correction pairing without manual annotation.

\paragraph{Incremental indexing.}
A manifest file (\texttt{index\_manifest.json}) tracks each file's GCS
generation number and content hash. Only files with changed content are
re-embedded, reducing re-indexing time from minutes to seconds.


% ─────────────────────────────────────────────────────────────────────────────
\subsubsection{Embedding Model Selection}
\label{subsec:embedding-selection}

Retrieval quality is bounded by the embedding model's discriminative power.
If semantically different queries map to similar vectors, no downstream
processing can recover the lost information---especially acute in math,
where minimal notational differences ($v(n)$ vs.\ $u(n)$) carry entirely
different semantics.

\paragraph{Initial experience.}
Our first implementation used a Google embedding model via Vertex AI.
During testing, it did not reliably distinguish between closely related
mathematical sequences---a $v(n)$ query would retrieve $u(n)$ chunks.
General-purpose models treat single-character differences as
insignificant~\cite{mansouri2019tangent, pfahler2020math_embeddings},
but in math they are decisive. This was a qualitative observation, not a
controlled experiment.

\paragraph{Migration to BGE-M3.}
We switched to BGE-M3~\cite{chen2024bge_m3}, built on XLM-RoBERTa with
8{,}192-token context. Its multilingual support (French + Arabic), long
context window, and multi-granularity training make it well-suited for our
heterogeneous corpus. The $v(n)$/$u(n)$ confusion no longer occurred after
migration.

The embedding function wraps the \texttt{BGEM3FlagModel} as a custom
ChromaDB \texttt{EmbeddingFunction}, loaded once as a singleton with
automatic GPU detection.


% ─────────────────────────────────────────────────────────────────────────────
\subsubsection{RAG Retrieval and Generation Engine}
\label{subsec:retrieval-engine}

The runtime pipeline is implemented in \texttt{rag\_engine.py}. It is
UI-agnostic: a single \texttt{.query(question, mode)} method returns a
structured result, callable from Streamlit, Jupyter, or CLI.

\paragraph{Two-stage retrieval.}
The first pass searches correction-type documents
(\texttt{is\_solution=true}), retrieving $k_1 = 10$ candidates. The
rationale is pedagogical: a previously corrected exercise provides not just
mathematical content but a \textit{rhetorical model}---the redaction style
the LLM can mimic.

If the best first-pass distance exceeds 1.2, a second pass retrieves up
to $k_2 = 8$ course-material documents for theorem backing. If both fail,
an unfiltered fallback search is used.

\paragraph{Similarity thresholds.}
Two L2 distance thresholds govern routing:
\begin{description}
    \item[Good match ($\leq 1.2$):] Strong correction match (Case~A). Only
    first-pass results are used.
    \item[Fallback ($\leq 1.6$):] Documents above this are discarded
    entirely.
\end{description}
These were calibrated empirically and are specific to BGE-M3 on our corpus.

\paragraph{Context compilation.}
Selected documents are formatted into XML-tagged context blocks with
metadata, role tags (\texttt{[\'ENONC\'E]} / \texttt{[CORRECTION]}), and
content truncated to 2{,}000 characters per document. Total context is
capped at 14{,}000 characters.

\paragraph{Prompt and generation.}
The system prompt defines the ``Monsieur Tounsi'' persona with:
anti-hallucination rules (use \textit{only} retrieved context), a mimicry
decision tree (Case~A: copy the correction style; Case~B: build from
theorems), a syllabus guard, the Derja sandwich convention, and
anti-injection rules. Generation uses Gemini ($T = 0.15$, max 4{,}096
tokens) with exponential-backoff retry.

\paragraph{Confidence and observability.}
Each query returns timing breakdowns, retrieval metadata, and a confidence
level: \textit{fort} (best distance $\leq 1.2$ and context $>$ 5{,}000
chars), \textit{moyen} (distance $\leq 1.6$), or \textit{faible} (no
useful results).


% ─────────────────────────────────────────────────────────────────────────────
\subsubsection{Paired Retrieval of Exercise and Correction}
\label{subsec:paired-retrieval}

During testing, we found that the system sometimes retrieved a correction
without its exercise statement, or vice versa. A correction is not
self-contained---the notation and constraints in the exercise define what
the correction solves.

\paragraph{Companion fetch mechanism.}
After two-stage retrieval selects its correction documents, the engine
performs a metadata lookup for each: using \texttt{chapter}, \texttt{type},
\texttt{year}, and \texttt{exo\_id}, it queries ChromaDB for chunks with
\texttt{is\_solution=false} (exercise statements). This uses ChromaDB's
direct metadata lookup, not vector search. Up to 3 companion chunks per
correction group are prepended to the context with an
\texttt{[\'ENONC\'E]} label.

This is a specialized form of parent-document
retrieval~\cite{gao2024rag_survey}, adapted to the paired structure of
mathematical exercises. To our knowledge, no prior work has implemented
exercise--correction companion retrieval in a mathematical tutoring system.


% ─────────────────────────────────────────────────────────────────────────────
\subsubsection{Connection to the Tunisian Bac Context}
\label{subsec:bac-context}

Four domain characteristics shaped the system design:

\begin{enumerate}
    \item \textbf{Curriculum boundaries.} The Bac covers a fixed set of ten
    chapters; examiners penalize out-of-program methods. Enforced through
    the syllabus guard and metadata-filtered retrieval.

    \item \textbf{Redaction conventions.} Bac corrections follow a strict
    rhetorical style. The two-stage retrieval prioritizes correction
    exemplars so the LLM can mimic this style.

    \item \textbf{Exercise--correction dependency.} Corrections are always
    read alongside their exercises. The companion fetch preserves this.

    \item \textbf{Bilingual context.} The corpus is in French with \LaTeX{}
    math; student queries may include Tunisian Derja. BGE-M3's multilingual
    training supports this.
\end{enumerate}
